\section{Introducción: el futuro de la IA agéntica}

Esta clase la presentan en conjunto dos invitados: \textbf{Chi Wang} (Microsoft Research, desarrollador principal del framework AutoGen) presenta los frameworks de IA agéntica y AutoGen; \textbf{Jerry Liu} (fundador y director general de LlamaIndex) presenta la construcción de asistentes de conocimiento multimodales.

\subsection{¿Qué es la IA agéntica?}

\begin{importantbox}{El valor central de la IA agéntica}
\begin{enumerate}
    \item \textbf{Interfaz en lenguaje natural}: el usuario describe sus necesidades en lenguaje natural y el agente las entiende y ejecuta automáticamente
    \item \textbf{Ejecución autónoma de tareas complejas}: mediante razonamiento de varios pasos y llamadas a herramientas, completa tareas con la mínima supervisión humana
    \item \textbf{Nueva arquitectura de software}: varios agentes colaboran para formar un sistema modular y componible de manera recursiva
\end{enumerate}
\end{importantbox}

El ejemplo muestra el proceso completo en que Magentic-One (basado en AutoGen) construye un sitio web de forma automática: analiza la tarea, instala las dependencias, varios agentes colaboran para escribir el código, lo compila y lo ejecuta, e incluso es capaz de autorrepararse después de que se elimina código clave.

\subsection{Resumen de la sección}

La IA agéntica no es solo una capa de aplicación de los LLM, sino que representa un paradigma completamente nuevo de desarrollo de software: pasar de «escribir código» a «orquestar agentes».

